% Discussion: the VVUQ tests, the quality and confidence of results against the
% external standards, the cost half of the thesis, comparison to prior work, and
% practical insights on real-life closeness. Built by the future pass from the
% instructions below. Do not process now. Figure 04 (binding matrix) and Figure 13
% (cost assessment) are placed here; \autoref both from the prose.

[Write the Discussion as an interpretive argument, not a recap. Carry comparisons
in left-aligned fixed-width tables and reference the figures with \autoref. The
through-line is the thesis: the assurance process, not generation, is the
decision-bearing work, and standards anchoring is what makes it defensible.]

[Subsection 1, what the VVUQ tests demonstrate. Discuss the tests performed in
\texttt{papers/VVUQ-02/codegen/tests/} (the 15 modules, indexed in
\texttt{papers/VVUQ-02/execution/01-foundation/test-suite.md}), arguing that the
test budget itself encodes the thesis: \texttt{test\_vvuq\_gates.py} carries 64 of
172 tests and \texttt{test\_safety.py} a further 12, so the assurance and safety
modules dominate the behavior modules. Connect specific test intents to the
standards they pin: the catastrophe-gate tests to IEC 80601-2-77 §201.x
hazardous-situation response and IEC 60601-1 single-fault safety, the balance tests
to ISO 13482 stability, the autonomy concordance tests to UL 4600 and IEEE 7009,
and the standards-corpus tests in \texttt{test\_standards\_corpus.py} to the
traceability of every binding. Cite \texttt{iec8060127}, \texttt{iec60601},
\texttt{iso13482}, \texttt{ul4600}, and \texttt{ieee7009}.]

[Subsection 2, quality and confidence against the external standards. This is the
credibility argument. Using the resolved per-gate standards table in
\texttt{papers/VVUQ-02/execution/03-vvuq/README.md} (\texttt{External standards
anchored here}) and the binding in
\texttt{papers/VVUQ-02/codegen/config/standards\_map.yaml} plus
\texttt{papers/VVUQ-02/inputs/standards/manifest.yaml}, argue that the decision
surface is credible because every gate decision is traceable to a published
designation rather than an ad hoc threshold. Map the three credibility questions to
their standards: verification and the 172-test floor to ASME V\&V 40-2018 §8 and
IEC 62304; validation against the independent references in
\texttt{papers/VVUQ-02/codegen/data/reference/} to ASME V\&V 40 §9; and uncertainty,
the coefficient-of-variation bounds and the ESCALATE path, to NASA-STD-7009A which
treats uncertainty as a first-class credibility factor. Tie the whole argument to
the FDA computational-model credibility framework. Introduce \autoref{fig:matrix}
here: the binding matrix shows that all fifteen governing standards are used and
that the catastrophe gates and the workhorse standard IEC 80601-2-77 carry the
densest anchoring. Cite \texttt{asmevv40}, \texttt{nasastd7009a},
\texttt{iec62304}, \texttt{iso14971}, and \texttt{fda2023credibility}.]

[Subsection 3, the cost and speed half of the thesis. From the \texttt{This
execution versus conventional high-end server processing} section of
\texttt{papers/VVUQ-02/execution/README.md}, discuss why the autonomous cloud run
was faster and less expensive than conventional verification: near-zero setup
(standard library plus four small packages with guarded imports), one integrated
agent loop (read, run, process, document, lint, commit, push with no tool context
switch), assurance breadth beyond the shipped tests in a single pass, and a
self-verifying record. Introduce \autoref{fig:cost} here as the explicit and honest
cost bridge, and state clearly that the index is illustrative and grounded in
direction only, because the repository records no dollar figures. Note the time
anchor that the 172-test suite runs in about 1.4 seconds. Then state honestly what
a conventional provisioned server still does better (live on-prem LLM intent and
judging, real Zenodo deposition, and physics-backed simulation), so the cost
argument does not overclaim.]

[Subsection 4, comparison to prior work and practical insights. Position the work
against autonomous clinical and code-generation systems: the deterministic
LLM-code-generation perspective of \texttt{trooskens2026compiled}, the autonomous
statistical-programming precision of \texttt{yan2026clinagent}, the validated
generative clinical pipeline of \texttt{carlisle2026pipeline}, and the digital-twin
VVUQ survey of \texttt{sel2025vvuq}; contrast each with this paper's emphasis on
assurance over generation for a single embodied surgeon. Then give practical,
honest insight on how close to real life this is: it is a planning and assurance
reference, not a robot controller, the platform is a clearly labelled hypothetical
2030 device, and every number is a simulation result, yet the framework checks
real steps off (determinism, traceable standards, pass-fail verification, recorded
human review, hand-back-to-human on escalation) so the industry does not have to
re-derive them. Frame the regulatory posture respectfully but opportunistically:
this is the kind of fast, standards-anchored assurance the FDA real-time clinical
trials direction invites (\texttt{fda2026realtime}), and US leadership in
surgical-humanoid VVUQ tied to external standards serves patient safety, efficacy,
and trial speed together.]

\begin{figure}[H]
   \centering
   \IfFileExists{Images/04-gate-standard-binding-matrix.png}{%
     \hspace{-0.95cm}%
     \includegraphics[width=1.1\linewidth, trim={0 0 0 0}, clip]{Images/04-gate-standard-binding-matrix.png}%
     % \vspace{-0.25cm}
   }{%
     \setlength{\fboxsep}{10pt}%
     \fbox{\parbox[c][6cm][c]{0.9\linewidth}{\centering\itshape Figure placeholder.
     The author adds \texttt{Images/04-gate-standard-binding-matrix.png}, rendered by
     \texttt{papers/VVUQ-02/imagegen/04-gate-standard-binding-matrix/04-gate-standard-binding-matrix.py}.}}%
   }
    \caption{Binding matrix of the ten VVUQ gates against the fifteen governing external standards, with every standard anchored by at least one gate.}
    \label{fig:matrix}
\end{figure}

\begin{figure}[H]
   \centering
   \IfFileExists{Images/13-assurance-cost-assessment.png}{%
     \hspace{-0.95cm}%
     \includegraphics[width=1.1\linewidth, trim={0 0 0 0}, clip]{Images/13-assurance-cost-assessment.png}%
     % \vspace{-0.25cm}
   }{%
     \setlength{\fboxsep}{10pt}%
     \fbox{\parbox[c][6cm][c]{0.9\linewidth}{\centering\itshape Figure placeholder.
     The author adds \texttt{Images/13-assurance-cost-assessment.png}, rendered by
     \texttt{papers/VVUQ-02/imagegen/13-assurance-cost-assessment/13-assurance-cost-assessment.py}.}}%
   }
    \caption{Assurance cost assessment, autonomous cloud VVUQ versus conventional verification, on an illustrative relative index grounded in direction.}
    \label{fig:cost}
\end{figure}
